\documentclass[../main.tex]{subfiles}
\begin{document}
\chapter*{GHI CHÚ VỀ VIỆC SỬ DỤNG CÔNG CỤ AI}
\addcontentsline{toc}{chapter}{GHI CHÚ VỀ VIỆC SỬ DỤNG CÔNG CỤ AI}
\thispagestyle{plain}

Trong quá trình thực hiện đồ án, tác giả có sử dụng công cụ trí tuệ nhân tạo (AI) như một
trợ lý soạn thảo văn bản và vẽ sơ đồ, dựa trên dữ liệu, mã nguồn và yêu cầu do tác giả cung
cấp. Ý tưởng, thiết kế hệ thống, mã nguồn, dữ liệu thực nghiệm và kết quả trong đồ án là do
tác giả tự thực hiện. Tác giả đã đọc, kiểm tra, chỉnh sửa và \textbf{chịu trách nhiệm về toàn
bộ nội dung} của quyển. Mức độ sử dụng AI cho từng hạng mục được liệt kê trong bảng dưới đây.

\vspace{0.6em}
\noindent\textbf{Công cụ đã sử dụng:}
\begin{itemize}[leftmargin=1.6em,topsep=2pt,itemsep=1pt]
  \item \textbf{Claude (Anthropic)} --- trợ lý AI hỗ trợ soạn thảo, biên tập văn bản và tham
  khảo khi viết mã nguồn.
  \item \textbf{PlantUML} --- công cụ sinh sơ đồ (kiến trúc, use case, tuần tự, quy trình) từ
  đặc tả do tác giả viết.
\end{itemize}

\vspace{0.4em}
\noindent\textbf{Quy ước mức độ:} \textit{Tự thực hiện} --- không sử dụng AI;
\textit{Tham khảo AI} --- tác giả tự làm, AI góp ý hoặc chỉnh sửa nhẹ;
\textit{Dùng AI} --- AI tạo ra, tác giả kiểm tra và chỉnh sửa lại.

\vspace{0.8em}
\begin{center}
\begin{tabular}{|p{9.6cm}|p{3.6cm}|}
\hline
\textbf{Hạng mục} & \textbf{Mức độ sử dụng AI} \\ \hline
Ý tưởng đề tài, thiết kế hệ thống, giải pháp chống gian lận & Tự thực hiện \\ \hline
Mã nguồn (Mini App, Cloud Functions, trang quản trị) & Tham khảo AI \\ \hline
Ảnh chụp màn hình ứng dụng & Tự thực hiện \\ \hline
Dữ liệu kiểm thử (72 ca), báo cáo bảo mật & Tự thực hiện \\ \hline
Nội dung lý thuyết chuyên môn (HMAC, FaceNet, mã QR) & Tham khảo AI \\ \hline
Bảng biểu, định dạng trình bày & Tham khảo AI \\ \hline
Văn xuôi thuyết minh các chương & Dùng AI \\ \hline
Sơ đồ (kiến trúc, use case, tuần tự, quy trình) & Dùng AI \\ \hline
\end{tabular}
\end{center}

\end{document}
